\documentclass[exercice,a4]{thedocument}%
\usepackage[french]{babel}%
\usepackage{tabularx}%
\usepackage{array}%
\startdatakeys
    \source{22GENMATMEAG1}
    \BO{2026}
    \cycle{4}
    \competencesDuSocle{CA,RA}
    \objectifsApprentissage{C4C2bc,C4A1aj}
    \motsCles{facteurs premiers,PGCD,probabilité}
\enddatakeys
\begin{document}%

\medskip

Une collectionneuse compte ses cartes Pokémon afin de les revendre. Elle possède 252 cartes de type \og feu \fg{} et 156 cartes de type \og terre \fg.
\begin{enumerate}
\item  \begin{enumerate}
\item  Parmi les trois propositions suivantes, laquelle correspond à la décomposition en produit de facteurs premiers du nombre 252 :

\begin{center}
\begin{tabularx}{9cm}{|*{3}{>{\centering \arraybackslash}X|}}\hline
Proposition 1 &Proposition 2& Proposition 3\\
 $2^2\times 9\times 7$ &$2\times 2\times 3\times 21$ &$2^2 \times 3^2\times  7$\\\hline
\end{tabularx}
\end{center}
\item Donner la décomposition en produit de facteurs premiers du nombre $156$.
\end{enumerate}
\item Elle veut réaliser des paquets identiques, c'est-à-dire contenant chacun le même nombre de cartes \og terre \fg{} et le même nombre de cartes \og feu \fg{} en utilisant toutes ses cartes.
\begin{enumerate}
\item Peut-elle faire 36 paquets ?
\item Quel est le nombre maximum de paquets qu'elle peut réaliser ?
\item Combien de cartes de chaque type contient alors chaque paquet ?
\end{enumerate}
\item Elle choisit une carte au hasard parmi toutes ses cartes. On suppose les cartes  indiscernables au toucher.

Calculer la probabilité que ce soit une carte de type \og terre \fg.
\end{enumerate}

\vspace{0.5cm}

\showthesource
\end{document}
